{
\captionsetup{type=figure}
\captionof{figure}{Trajectory for instance \texttt{sympy\_\_sympy-14774}.}
\begin{issuebox}[sympy\_\_sympy-14774]We're currently solving the following issue within our repository. Here's the issue text:

ISSUE:

Latex printer does not support full inverse trig function names for acsc and asec

For example

`latex(asin(x), inv\_trig\_style="full")' works as expected returning `'\\arcsin{\\left (x \\right )}''

But `latex(acsc(x), inv\_trig\_style="full")' gives `'\\operatorname{acsc}{\\left (x \\right )}'' instead of `'\\operatorname{arccsc}{\\left (x \\right )}''



A fix seems to be to change line 743 of sympy/printing/latex.py from

`inv\_trig\_table = ["asin", "acos", "atan", "acot"]' to

`inv\_trig\_table = ["asin", "acos", "atan", "acsc", "asec", "acot"]'









INSTRUCTIONS: (31 lines)

\begin{CodeVerbatim}[breaklines=true, breakanywhere=true]
(Open file: n/a)
(Current directory: /sympy__sympy)
bash-\$
\end{CodeVerbatim}
\end{issuebox}
\begin{agentbox}[sympy\_\_sympy-14774]
First, I'll create a Python script to reproduce the issue as described. This script will use the `latex` function from SymPy to print the inverse trigonometric functions with the `inv\_trig\_style=''full''` option. I'll name this script `reproduce\_issue.py`.



\hrulefill
\begin{CodeVerbatim}[breaklines=true, breakanywhere=true]
create reproduce_issue.py
\end{CodeVerbatim}
\end{agentbox}
\begin{observationbox}[sympy\_\_sympy-14774]\begin{Verbatim}[breaklines=true, breakanywhere=true]
[File: /sympy__sympy/reproduce_issue.py (1 lines total)]
1:

(Open file: n/a)
(Current directory: /sympy__sympy)
bash-$
\end{Verbatim}
\end{observationbox}
\begin{agentbox}[sympy\_\_sympy-14774]
Now I will input the code provided in the issue into the `reproduce\_issue.py` file to replicate the bug.



\hrulefill
\begin{CodeVerbatim}[breaklines=true, breakanywhere=true]
edit 1:1
from sympy import symbols, latex, asin, acsc

x = symbols('x')
print(latex(asin(x), inv_trig_style="full"))
print(latex(acsc(x), inv_trig_style="full"))
end_of_edit
\end{CodeVerbatim}
\end{agentbox}
\begin{observationbox}[sympy\_\_sympy-14774]\begin{Verbatim}[breaklines=true, breakanywhere=true]
[File: /sympy__sympy/reproduce_issue.py (5 lines total)]
1:from sympy import symbols, latex, asin, acsc
2:
3:x = symbols('x')
4:print(latex(asin(x), inv_trig_style="full"))
5:print(latex(acsc(x), inv_trig_style="full"))
File updated. Please review the changes and make sure they are correct (correct indentation, no duplicate lines, etc). Edit the file again if necessary.

(Open file: /sympy__sympy/reproduce_issue.py)
(Current directory: /sympy__sympy)
bash-$
\end{Verbatim}
\end{observationbox}
\begin{agentbox}[sympy\_\_sympy-14774]
The script has been updated with the code to reproduce the issue. Now, let's run the script to confirm that the issue exists.



\hrulefill
\begin{CodeVerbatim}[breaklines=true, breakanywhere=true]
python reproduce_issue.py
\end{CodeVerbatim}
\end{agentbox}
\begin{observationbox}[sympy\_\_sympy-14774]\begin{Verbatim}[breaklines=true, breakanywhere=true]
/sympy__sympy/sympy/core/basic.py:3: DeprecationWarning: Using or importing the ABCs from 'collections' instead of from 'collections.abc' is deprecated since Python 3.3, and in 3.10 it will stop working
  from collections import Mapping, defaultdict
/sympy__sympy/sympy/core/expr.py:12: DeprecationWarning: Using or importing the ABCs from 'collections' instead of from 'collections.abc' is deprecated since Python 3.3, and in 3.10 it will stop working
  from collections import defaultdict, Iterable
/sympy__sympy/sympy/core/containers.py:271: DeprecationWarning: Using or importing the ABCs from 'collections' instead of from 'collections.abc' is deprecated since Python 3.3, and in 3.10 it will stop working
  class OrderedSet(collections.MutableSet):
/sympy__sympy/sympy/solvers/diophantine.py:3188: SyntaxWarning: "is" with a literal. Did you mean "=="?
  if feasible is 1:  # it's prime and k == 2
/sympy__sympy/sympy/plotting/plot.py:28: DeprecationWarning: Using or importing the ABCs from 'collections' instead of from 'collections.abc' is deprecated since Python 3.3, and in 3.10 it will stop working
  from collections import Callable
\arcsin{\left (x \right )}
\operatorname{acsc}{\left (x \right )}

(Open file: /sympy__sympy/reproduce_issue.py)
(Current directory: /sympy__sympy)
bash-$
\end{Verbatim}
\end{observationbox}
\begin{agentbox}[sympy\_\_sympy-14774]
The output of the script confirms the issue: the `latex` function returns `'{\textbackslash}arcsin\{{\textbackslash}left (x {\textbackslash}right )\}'` for `asin(x)` as expected, but for `acsc(x)` it returns `'{\textbackslash}operatorname\{acsc\}\{{\textbackslash}left (x {\textbackslash}right )\}'` instead of `'{\textbackslash}operatorname\{arccsc\}\{{\textbackslash}left (x {\textbackslash}right )\}'`. This matches the behavior described in the issue.



Now that we have confirmed the issue, we can proceed to fix it. According to the issue, we need to modify line 743 of `sympy/printing/latex.py`. Let's find the `latex.py` file in the repository.



\hrulefill
\begin{CodeVerbatim}[breaklines=true, breakanywhere=true]
find_file "latex.py" sympy/printing
\end{CodeVerbatim}
\end{agentbox}
\begin{observationbox}[sympy\_\_sympy-14774]\begin{Verbatim}[breaklines=true, breakanywhere=true]
Found 1 matches for "latex.py" in /sympy__sympy/sympy/printing:
/sympy__sympy/sympy/printing/latex.py

(Open file: /sympy__sympy/reproduce_issue.py)
(Current directory: /sympy__sympy)
bash-$
\end{Verbatim}
\end{observationbox}
\begin{agentbox}[sympy\_\_sympy-14774]
The `latex.py` file has been located in the `sympy/printing` directory. Let's open the file and navigate to line 743 to make the necessary changes.



\hrulefill
\begin{CodeVerbatim}[breaklines=true, breakanywhere=true]
open sympy/printing/latex.py 743
\end{CodeVerbatim}
\end{agentbox}
\begin{observationbox}[sympy\_\_sympy-14774]\begin{Verbatim}[breaklines=true, breakanywhere=true]
[File: /sympy__sympy/sympy/printing/latex.py (2350 lines total)]
(727 more lines above)
728:        if hasattr(self, '_print_' + func) and \
729:            not isinstance(expr.func, UndefinedFunction):
730:            return getattr(self, '_print_' + func)(expr, exp)
731:        else:
732:            args = [ str(self._print(arg)) for arg in expr.args ]
733:            # How inverse trig functions should be displayed, formats are:
734:            # abbreviated: asin, full: arcsin, power: sin^-1
735:            inv_trig_style = self._settings['inv_trig_style']
736:            # If we are dealing with a power-style inverse trig function
737:            inv_trig_power_case = False
738:            # If it is applicable to fold the argument brackets
739:            can_fold_brackets = self._settings['fold_func_brackets'] and \
740:                len(args) == 1 and \
741:                not self._needs_function_brackets(expr.args[0])
742:
743:            inv_trig_table = ["asin", "acos", "atan", "acot"]
744:
745:            # If the function is an inverse trig function, handle the style
746:            if func in inv_trig_table:
747:                if inv_trig_style == "abbreviated":
748:                    func = func
749:                elif inv_trig_style == "full":
750:                    func = "arc" + func[1:]
751:                elif inv_trig_style == "power":
752:                    func = func[1:]
753:                    inv_trig_power_case = True
754:
755:                    # Can never fold brackets if we're raised to a power
756:                    if exp is not None:
757:                        can_fold_brackets = False
758:
759:            if inv_trig_power_case:
760:                if func in accepted_latex_functions:
761:                    name = r"\%s^{-1}" % func
762:                else:
763:                    name = r"\operatorname{%s}^{-1}" % func
764:            elif exp is not None:
765:                name = r'%s^{%s}' % (self._hprint_Function(func), exp)
766:            else:
767:                name = self._hprint_Function(func)
768:
769:            if can_fold_brackets:
770:                if func in accepted_latex_functions:
771:                    # Wrap argument safely to avoid parse-time conflicts
772:                    # with the function name itself
773:                    name += r" {%s}"
774:                else:
775:                    name += r"%s"
776:            else:
777:                name += r"{\left (%s \right )}"
778:
779:            if inv_trig_power_case and exp is not None:
780:                name += r"^{%s}" % exp
781:
782:            return name % ",".join(args)
783:
784:    def _print_UndefinedFunction(self, expr):
785:        return self._hprint_Function(str(expr))
786:
787:    @property
788:    def _special_function_classes(self):
789:        from sympy.functions.special.tensor_functions import KroneckerDelta
790:        from sympy.functions.special.gamma_functions import gamma, lowergamma
791:        from sympy.functions.special.beta_functions import beta
792:        from sympy.functions.special.delta_functions import DiracDelta
793:        from sympy.functions.special.error_functions import Chi
794:        return {KroneckerDelta: r'\delta',
795:                gamma:  r'\Gamma',
796:                lowergamma: r'\gamma',
797:                beta: r'\operatorname{B}',
798:                DiracDelta: r'\delta',
799:                Chi: r'\operatorname{Chi}'}
800:
801:    def _print_FunctionClass(self, expr):
802:        for cls in self._special_function_classes:
803:            if issubclass(expr, cls) and expr.__name__ == cls.__name__:
804:                return self._special_function_classes[cls]
805:        return self._hprint_Function(str(expr))
806:
807:    def _print_Lambda(self, expr):
808:        symbols, expr = expr.args
809:
810:        if len(symbols) == 1:
811:            symbols = self._print(symbols[0])
812:        else:
813:            symbols = self._print(tuple(symbols))
814:
815:        args = (symbols, self._print(expr))
816:        tex = r"\left( %s \mapsto %s \right)" % (symbols, self._print(expr))
817:
818:        return tex
819:
820:    def _print_Min(self, expr, exp=None):
821:        args = sorted(expr.args, key=default_sort_key)
822:        texargs = [r"%s" % self._print(symbol) for symbol in args]
823:        tex = r"\min\left(%s\right)" % ", ".join(texargs)
824:
825:        if exp is not None:
826:            return r"%s^{%s}" % (tex, exp)
827:        else:
(1523 more lines below)

(Open file: /sympy__sympy/reproduce_issue.py)
(Current directory: /sympy__sympy)
bash-$
\end{Verbatim}
\end{observationbox}
\begin{agentbox}[sympy\_\_sympy-14774]
We have the `latex.py` file open at the relevant section. The issue suggests that we need to add `''acsc''` and `''asec''` to the `inv\_trig\_table` list on line 743. Let's make the necessary edit to the code.



\hrulefill
\begin{CodeVerbatim}[breaklines=true, breakanywhere=true]
edit 743:743
inv_trig_table = ["asin", "acos", "atan", "acsc", "asec", "acot"]
end_of_edit
\end{CodeVerbatim}
\end{agentbox}
\begin{observationbox}[sympy\_\_sympy-14774]\begin{Verbatim}[breaklines=true, breakanywhere=true]
Your proposed edit has introduced new syntax error(s). Please understand the fixes and retry your edit commmand.

ERRORS:
- E999 IndentationError: unexpected indent

This is how your edit would have looked if applied
-------------------------------------------------
[File: /sympy__sympy/sympy/printing/latex.py (2350 lines total)]
(736 more lines above)
737:            inv_trig_power_case = False
738:            # If it is applicable to fold the argument brackets
739:            can_fold_brackets = self._settings['fold_func_brackets'] and \
740:                len(args) == 1 and \
741:                not self._needs_function_brackets(expr.args[0])
742:
743:inv_trig_table = ["asin", "acos", "atan", "acsc", "asec", "acot"]
744:
745:            # If the function is an inverse trig function, handle the style
746:            if func in inv_trig_table:
747:                if inv_trig_style == "abbreviated":
(1603 more lines below)
-------------------------------------------------

This is the original code before your edit
-------------------------------------------------
[File: /sympy__sympy/sympy/printing/latex.py (2350 lines total)]
(737 more lines above)
738:            # If it is applicable to fold the argument brackets
739:            can_fold_brackets = self._settings['fold_func_brackets'] and \
740:                len(args) == 1 and \
741:                not self._needs_function_brackets(expr.args[0])
742:
743:            inv_trig_table = ["asin", "acos", "atan", "acot"]
744:
745:            # If the function is an inverse trig function, handle the style
746:            if func in inv_trig_table:
747:                if inv_trig_style == "abbreviated":
748:                    func = func
(1602 more lines below)
-------------------------------------------------
Your changes have NOT been applied. Please fix your edit command and try again.
You either need to 1) Specify the correct start/end line arguments or 2) Correct your edit code.
DO NOT re-run the same failed edit command. Running it again will lead to the same error.

(Open file: /sympy__sympy/sympy/printing/latex.py)
(Current directory: /sympy__sympy)
bash-$
\end{Verbatim}
\end{observationbox}
\begin{agentbox}[sympy\_\_sympy-14774]
The edit command failed due to incorrect indentation. I will correct the indentation and try the edit command again.



\hrulefill
\begin{CodeVerbatim}[breaklines=true, breakanywhere=true]
edit 743:743
            inv_trig_table = ["asin", "acos", "atan", "acsc", "asec", "acot"]
end_of_edit
\end{CodeVerbatim}
\end{agentbox}
\begin{observationbox}[sympy\_\_sympy-14774]\begin{Verbatim}[breaklines=true, breakanywhere=true]
[File: /sympy__sympy/sympy/printing/latex.py (2350 lines total)]
(692 more lines above)
693:        if isinstance(e, AssocOp):
694:            return r"%s\left(%s\right)" % (tex, self._print(e))
695:        else:
696:            return r"%s %s" % (tex, self._print(e))
697:
698:    def _hprint_Function(self, func):
699:        r'''
700:        Logic to decide how to render a function to latex
701:          - if it is a recognized latex name, use the appropriate latex command
702:          - if it is a single letter, just use that letter
703:          - if it is a longer name, then put \operatorname{} around it and be
704:            mindful of undercores in the name
705:        '''
706:        func = self._deal_with_super_sub(func)
707:        if func in accepted_latex_functions:
708:            name = r"\%s" % func
709:        elif len(func) == 1 or func.startswith('\\'):
710:            name = func
711:        else:
712:            name = r"\operatorname{%s}" % func
713:        return name
714:
715:    def _print_Function(self, expr, exp=None):
716:        r'''
717:        Render functions to LaTeX, handling functions that LaTeX knows about
718:        e.g., sin, cos, ... by using the proper LaTeX command (\sin, \cos, ...).
719:        For single-letter function names, render them as regular LaTeX math
720:        symbols. For multi-letter function names that LaTeX does not know
721:        about, (e.g., Li, sech) use \operatorname{} so that the function name
722:        is rendered in Roman font and LaTeX handles spacing properly.
723:
724:        expr is the expression involving the function
725:        exp is an exponent
726:        '''
727:        func = expr.func.__name__
728:        if hasattr(self, '_print_' + func) and \
729:            not isinstance(expr.func, UndefinedFunction):
730:            return getattr(self, '_print_' + func)(expr, exp)
731:        else:
732:            args = [ str(self._print(arg)) for arg in expr.args ]
733:            # How inverse trig functions should be displayed, formats are:
734:            # abbreviated: asin, full: arcsin, power: sin^-1
735:            inv_trig_style = self._settings['inv_trig_style']
736:            # If we are dealing with a power-style inverse trig function
737:            inv_trig_power_case = False
738:            # If it is applicable to fold the argument brackets
739:            can_fold_brackets = self._settings['fold_func_brackets'] and \
740:                len(args) == 1 and \
741:                not self._needs_function_brackets(expr.args[0])
742:
743:            inv_trig_table = ["asin", "acos", "atan", "acsc", "asec", "acot"]
744:
745:            # If the function is an inverse trig function, handle the style
746:            if func in inv_trig_table:
747:                if inv_trig_style == "abbreviated":
748:                    func = func
749:                elif inv_trig_style == "full":
750:                    func = "arc" + func[1:]
751:                elif inv_trig_style == "power":
752:                    func = func[1:]
753:                    inv_trig_power_case = True
754:
755:                    # Can never fold brackets if we're raised to a power
756:                    if exp is not None:
757:                        can_fold_brackets = False
758:
759:            if inv_trig_power_case:
760:                if func in accepted_latex_functions:
761:                    name = r"\%s^{-1}" % func
762:                else:
763:                    name = r"\operatorname{%s}^{-1}" % func
764:            elif exp is not None:
765:                name = r'%s^{%s}' % (self._hprint_Function(func), exp)
766:            else:
767:                name = self._hprint_Function(func)
768:
769:            if can_fold_brackets:
770:                if func in accepted_latex_functions:
771:                    # Wrap argument safely to avoid parse-time conflicts
772:                    # with the function name itself
773:                    name += r" {%s}"
774:                else:
775:                    name += r"%s"
776:            else:
777:                name += r"{\left (%s \right )}"
778:
779:            if inv_trig_power_case and exp is not None:
780:                name += r"^{%s}" % exp
781:
782:            return name % ",".join(args)
783:
784:    def _print_UndefinedFunction(self, expr):
785:        return self._hprint_Function(str(expr))
786:
787:    @property
788:    def _special_function_classes(self):
789:        from sympy.functions.special.tensor_functions import KroneckerDelta
790:        from sympy.functions.special.gamma_functions import gamma, lowergamma
791:        from sympy.functions.special.beta_functions import beta
792:        from sympy.functions.special.delta_functions import DiracDelta
(1558 more lines below)
File updated. Please review the changes and make sure they are correct (correct indentation, no duplicate lines, etc). Edit the file again if necessary.

(Open file: /sympy__sympy/sympy/printing/latex.py)
(Current directory: /sympy__sympy)
bash-$
\end{Verbatim}
\end{observationbox}
\begin{agentbox}[sympy\_\_sympy-14774]
The `latex.py` file has been successfully updated with the correct entries in the `inv\_trig\_table`. Now, to verify that the fix works, we should re-run the `reproduce\_issue.py` script to see if the output for `acsc(x)` is now `'{\textbackslash}operatorname\{arccsc\}\{{\textbackslash}left (x {\textbackslash}right )\}'`.



\hrulefill
\begin{CodeVerbatim}[breaklines=true, breakanywhere=true]
python reproduce_issue.py
\end{CodeVerbatim}
\end{agentbox}
\begin{observationbox}[sympy\_\_sympy-14774]\begin{Verbatim}[breaklines=true, breakanywhere=true]
/sympy__sympy/sympy/core/basic.py:3: DeprecationWarning: Using or importing the ABCs from 'collections' instead of from 'collections.abc' is deprecated since Python 3.3, and in 3.10 it will stop working
  from collections import Mapping, defaultdict
/sympy__sympy/sympy/core/expr.py:12: DeprecationWarning: Using or importing the ABCs from 'collections' instead of from 'collections.abc' is deprecated since Python 3.3, and in 3.10 it will stop working
  from collections import defaultdict, Iterable
/sympy__sympy/sympy/core/containers.py:271: DeprecationWarning: Using or importing the ABCs from 'collections' instead of from 'collections.abc' is deprecated since Python 3.3, and in 3.10 it will stop working
  class OrderedSet(collections.MutableSet):
/sympy__sympy/sympy/plotting/plot.py:28: DeprecationWarning: Using or importing the ABCs from 'collections' instead of from 'collections.abc' is deprecated since Python 3.3, and in 3.10 it will stop working
  from collections import Callable
\arcsin{\left (x \right )}
\operatorname{arccsc}{\left (x \right )}

(Open file: /sympy__sympy/sympy/printing/latex.py)
(Current directory: /sympy__sympy)
bash-$
\end{Verbatim}
\end{observationbox}
\begin{agentbox}[sympy\_\_sympy-14774]
The output for `acsc(x)` is now correctly formatted as `'{\textbackslash}operatorname\{arccsc\}\{{\textbackslash}left (x {\textbackslash}right )\}'`, which means the fix has been successful. The deprecation warnings are unrelated to the issue at hand and are a separate concern that would need to be addressed by updating the usage of the collections module throughout the SymPy codebase.



Since the fix for the issue is confirmed, we can now remove the `reproduce\_issue.py` script as it is no longer needed.



\hrulefill
\begin{CodeVerbatim}[breaklines=true, breakanywhere=true]
rm reproduce_issue.py
\end{CodeVerbatim}
\end{agentbox}
\begin{observationbox}[sympy\_\_sympy-14774]\begin{Verbatim}[breaklines=true, breakanywhere=true]

(Open file: /sympy__sympy/sympy/printing/latex.py)
(Current directory: /sympy__sympy)
bash-$
\end{Verbatim}
\end{observationbox}
\begin{agentbox}[sympy\_\_sympy-14774]
The `reproduce\_issue.py` script has been successfully removed. Now that the issue has been fixed and verified, we can submit the changes to the code base. Let's proceed with the submission using the `submit` command.



\hrulefill
\begin{CodeVerbatim}[breaklines=true, breakanywhere=true]
submit
\end{CodeVerbatim}
\end{agentbox}
\begin{observationbox}[sympy\_\_sympy-14774]\begin{CodeVerbatim}[breaklines=true, breakanywhere=true]

diff --git a/sympy/printing/latex.py b/sympy/printing/latex.py
--- a/sympy/printing/latex.py
+++ b/sympy/printing/latex.py
@@ -740,7 +740,7 @@ def _print_Function(self, expr, exp=None):
                 len(args) == 1 and \
                 not self._needs_function_brackets(expr.args[0])
 
-            inv_trig_table = ["asin", "acos", "atan", "acot"]
+            inv_trig_table = ["asin", "acos", "atan", "acsc", "asec", "acot"]
 
             # If the function is an inverse trig function, handle the style
             if func in inv_trig_table:
\end{CodeVerbatim}
\end{observationbox}

\begin{goldpatchbox}[sympy\_\_sympy-14774]
\begin{CodeVerbatim}[breaklines=true, breakanywhere=true]
diff --git a/sympy/printing/latex.py b/sympy/printing/latex.py
--- a/sympy/printing/latex.py
+++ b/sympy/printing/latex.py
@@ -740,7 +740,7 @@ def _print_Function(self, expr, exp=None):
                 len(args) == 1 and \
                 not self._needs_function_brackets(expr.args[0])
 
-            inv_trig_table = ["asin", "acos", "atan", "acot"]
+            inv_trig_table = ["asin", "acos", "atan", "acsc", "asec", "acot"]
 
             # If the function is an inverse trig function, handle the style
             if func in inv_trig_table:
\end{CodeVerbatim}
\end{goldpatchbox}
}